%%
%% part4-advanced/ch23-modules.tex
%%
%% Chapter 23: Modular Structure — The 19 modules of
%% MatinBook, their roles, and the loading order.
%%

\chapter{ساختار ماژولار}
\label{chap:modules}

\section{چرا این موضوع مهم است؟}
\label{sec:modules-why}

\lr{MatinBook} یک کلاس ماژولار است. یعنی
به‌جای یک فایل بزرگ، از ۱۹ ماژول مستقل
تشکیل شده که هر یک مسئولیت مشخصی دارند.
این معماری، نگهداری، توسعه و شخصی‌سازی
کلاس را بسیار آسان‌تر می‌کند.

درک ساختار ماژولار، برای شخصی‌سازی پیشرفته
و افزودن قابلیت‌های جدید ضروری است. در این
فصل، با این ساختار آشنا می‌شوید.

\mnote{معماری ماژولار یعنی اگر فقط می‌خواهید یک ماژول را تغییر دهید، نیازی به دست زدن به بقیه‌ی کلاس ندارید.}

\section{فهرست ماژول‌ها}
\label{sec:modules-list}

\lr{MatinBook} ۱۹ ماژول دارد که در سه دسته
تقسیم می‌شوند:

\subsection{هسته (۲ ماژول)}
\label{sec:modules-core}

\begin{itemize}
    \item \file{mb-core.sty}: بارگذارنده‌ی
    بسته‌های پایه.
    \item \file{mb-utils.sty}: دستورات کمکی.
\end{itemize}

\subsection{ماژول‌های کاربردی (۱۰ ماژول)}
\label{sec:modules-feature}

\begin{itemize}
    \item \file{mb-math.sty}: ریاضیات.
    \item \file{mb-boxes.sty}: استایل‌های
    جعبه.
    \item \file{mb-theorem.sty}: محیط‌های
    قضیه.
    \item \file{mb-code.sty}: کد.
    \item \file{mb-layout.sty}: چیدمان.
    \item \file{mb-headings.sty}: عناوین.
    \item \file{mb-references.sty}: ارجاعات.
    \item \file{mb-graphics.sty}: گرافیک.
    \item \file{mb-index.sty}: نمایه.
    \item \file{mb-table.sty}: جداول.
\end{itemize}

\subsection{تم و لوکال (۷ ماژول)}
\label{sec:modules-theme}

\begin{itemize}
    \item \file{mb-theme-colors.sty}: پالت
    رنگی.
    \item \file{mb-theme-cover.sty}: جلد.
    \item \file{mb-theme-default.sty}:
    بارگذارنده‌ی تم.
    \item \file{mb-typography.sty}:
    تایپوگرافی.
    \item \file{fa-IR.sty}: لوکال فارسی.
    \item \file{en-US.sty}: لوکال انگلیسی.
\end{itemize}

\mnote{جمع این ماژول‌ها ۱۹ عدد است: ۲ هسته + ۱۰ کاربردی + ۳ تم + ۲ لوکال + ۱ تایپوگرافی + ۱ کلاس اصلی.}

\section{نقش هر ماژول}
\label{sec:modules-roles}

جدول \cref{tab:module-roles} نقش هر ماژول
را خلاصه می‌کند.

\begin{matintable}{نقش ماژول‌های \lr{MatinBook}}{tab:module-roles}
\begin{tblr}{
    width = \textwidth,
    colspec = {l X[l]},
    hlines, vlines,
    colsep = 8pt,
    rowsep = 4pt,
    row{1} = {font=\bfseries},
}
ماژول & نقش \\
\lr{mb-core} & بارگذاری بسته‌های پایه \\
\lr{mb-utils} & دستورات کمکی \\
\lr{mb-math} & عملگرها و جداکننده‌های ریاضی \\
\lr{mb-boxes} & استایل‌های \lr{tcolorbox} \\
\lr{mb-theorem} & محیط‌های قضیه \\
\lr{mb-code} & \lr{minted} و پشتیبانی فارسی \\
\lr{mb-layout} & هندسه، سرصفحه، پاصفحه \\
\lr{mb-headings} & سبک عناوین \\
\lr{mb-references} & \lr{hyperref} و \lr{cleveref} \\
\lr{mb-graphics} & \lr{TikZ} و \lr{PGFPlots} \\
\lr{mb-index} & \lr{xindy} \\
\lr{mb-table} & \lr{tabularray} \\
\lr{mb-theme-colors} & پالت رنگی \lr{CMYK} \\
\lr{mb-theme-cover} & جلد جلو و عقب \\
\lr{mb-theme-default} & بارگذارنده‌ی تم \\
\lr{mb-typography} & \lr{microtype} و کشیده \\
\lr{fa-IR} & لوکال فارسی \\
\lr{en-US} & لوکال انگلیسی \\
\lr{matinbook.cls} & کلاس اصلی \\
\end{tblr}
\end{matintable}

\section{ترتیب بارگذاری}
\label{sec:modules-order}

ترتیب بارگذاری ماژول‌ها \textbf{حیاتی}
است. \lr{MatinBook} از سه فاز استفاده می‌کند:

\subsection{فاز ۱ — قبل از \lr{xepersian}}
\label{sec:modules-phase1}

\begin{latin}
\begin{minted}{text}
1.  mb-core
2.  mb-theme-colors
3.  mb-table
4.  mb-math
5.  mb-boxes
6.  mb-theorem
7.  mb-code
8.  mb-layout
9.  mb-headings
10. mb-references
11. mb-graphics
12. mb-index
13. mb-utils
14. mb-theme-cover
\end{minted}
\end{latin}

\mnote{در فاز ۱، ترتیب دقیق مهم است: \lr{mb-theme-colors} باید قبل از مصرف‌کنندگان (مثل \lr{mb-boxes}) بارگذاری شود.}

\subsection{فاز ۲ — \lr{xepersian}}
\label{sec:modules-phase2}

\begin{latin}
\begin{minted}{text}
15. xepersian
\end{minted}
\end{latin}

\mnote{\lr{xepersian} باید \textbf{آخرین} بسته‌ی فاز ۲ باشد. اگر بسته‌ای بعد از آن بارگذاری شود، تعاریف \lr{bidi} را بازنویسی می‌کند.}

\subsection{فاز ۳ — بعد از \lr{xepersian}}
\label{sec:modules-phase3}

\begin{latin}
\begin{minted}{text}
16. mb-typography
17. fa-IR
18. mb-theme-default
\end{minted}
\end{latin}

\mnote{در فاز ۳، فونت‌ها، تایپوگرافی و لوکال بارگذاری می‌شوند. این‌ها نیاز به \lr{xepersian} دارند.}

\section{چرا این ترتیب مهم است؟}
\label{sec:modules-why-order}

\begin{enumerate}
    \item \textbf{\lr{xepersian} آخرین بسته}:
    چون \lr{bidi} و سایر بسته‌ها را بازنویسی
    می‌کند. اگر بسته‌ای بعد از آن بارگذاری
    شود، تعاریف آن از بین می‌رود.

    \item \textbf{رنگ‌ها قبل از مصرف‌کنندگان}:
    \file{mb-theme-colors} قبل از
    \file{mb-boxes}، \file{mb-code}،
    \file{mb-layout}، \file{mb-headings} و
    \file{mb-graphics} بارگذاری می‌شود.

    \item \textbf{فونت‌ها بعد از \lr{xepersian}}:
    \cmd{settextfont} توسط \lr{xepersian}
    تعریف می‌شود، پس فونت‌ها باید در فاز ۳
    تنظیم شوند.
\end{enumerate}

\mnote{تغییر این ترتیب، باعث خطاهای عجیب و غریب می‌شود. اگر ماژول جدیدی اضافه می‌کنید، حتماً آن را در فاز مناسب قرار دهید.}

\section{ساختار پوشه‌ها}
\label{sec:modules-structure}

ماژول‌ها در پوشه‌ی \file{tex/} قرار دارند:

\begin{latin}
\begin{minted}{text}
tex/
├── core/
│   ├── mb-core.sty
│   └── mb-utils.sty
├── locales/
│   ├── fa-IR.sty
│   └── en-US.sty
├── modules/
│   ├── boxes/mb-boxes.sty
│   ├── code/mb-code.sty
│   ├── graphics/mb-graphics.sty
│   ├── index/mb-index.sty
│   ├── layout/
│   │   ├── mb-layout.sty
│   │   └── mb-headings.sty
│   ├── math/mb-math.sty
│   ├── references/mb-references.sty
│   ├── table/mb-table.sty
│   ├── theorem/mb-theorem.sty
│   └── typography/mb-typography.sty
└── themes/
    └── default/
        ├── mb-theme-colors.sty
        ├── mb-theme-cover.sty
        └── mb-theme-default.sty
\end{minted}
\end{latin}

\mnote{هر ماژول در پوشه‌ی مخصوص خود قرار دارد. این ساختار، پیدا کردن ماژول‌ها را آسان می‌کند.}

\section{افزودن ماژول جدید}
\label{sec:modules-adding}

برای افزودن یک ماژول جدید، این مراحل را
دنبال کنید:

\begin{enumerate}
    \item یک پوشه‌ی جدید در
    \file{tex/modules/} بسازید (مثلاً
    \file{mymodule/}).

    \item یک فایل \file{.sty} در آن بسازید
    (مثلاً \file{mb-mymodule.sty}).

    \item فایل را با ساختار استاندارد
    بنویسید.

    \item در \file{matinbook.cls}، ماژول را
    در فاز مناسب بارگذاری کنید.
\end{enumerate}

\subsection{ساختار یک ماژول}
\label{sec:modules-template}

\begin{latin}
\begin{minted}{latex}
\NeedsTeXFormat{LaTeX2e}[2023/06/01]
\ProvidesPackage{mb-mymodule}[2026/10/08 v1.2 MatinBook My Module]

%========================================================================
% MB-MYMODULE — DESCRIPTION
%========================================================================

% |\pc{کد ماژول...}|%

\PackageInfo{mb-mymodule}{My module loaded}

\endinput
\end{minted}
\end{latin}

\mnote{هر ماژول باید \lr{\textbackslash NeedsTeXFormat}، \lr{\textbackslash ProvidesPackage}، \lr{\textbackslash PackageInfo} و \lr{\textbackslash endinput} داشته باشد.}

\section{ساختار کلاس اصلی}
\label{sec:modules-cls}

فایل \file{matinbook.cls} سه فاز دارد:

\begin{latin}
\begin{minted}{latex}
% |\pc{فاز ۱: بسته‌های پایه}|%
\RequirePackage{mb-core}
\RequirePackage{mb-theme-colors}
\RequirePackage{mb-table}
\RequirePackage{mb-math}
\RequirePackage{mb-boxes}
\RequirePackage{mb-theorem}
\RequirePackage{mb-code}
\RequirePackage{mb-layout}
\RequirePackage{mb-headings}
\RequirePackage{mb-references}
\RequirePackage{mb-graphics}
\RequirePackage{mb-index}
\RequirePackage{mb-utils}
\RequirePackage{mb-theme-cover}

% |\pc{فاز ۲:}|\lr{xepersian}|%
\RequirePackage{xepersian}

% |\pc{فاز ۳: تنظیمات فارسی}|%
% |\pc{فونت‌ها...}|%
\RequirePackage{mb-typography}
\RequirePackage{fa-IR}
\end{minted}
\end{latin}

\mnote{در \lr{matinbook.cls}، هیچ ماژولی نباید خارج از این سه فاز بارگذاری شود.}

\section{اشکال‌زدایی}
\label{sec:modules-debugging}

اگر خطایی رخ داد، این مراحل را انجام
دهید:

\begin{enumerate}
    \item \textbf{بررسی فاز:} آیا ماژول در
    فاز درست بارگذاری شده؟

    \item \textbf{بررسی ترتیب:} آیا ترتیب
    ماژول‌ها درست است؟

    \item \textbf{بررسی وابستگی:} آیا ماژول
    به بسته‌ای نیاز دارد که بارگذاری
    نشده؟

    \item \textbf{بررسی نسخه:} آیا نسخه‌ی
    \lr{LaTeX} و بسته‌ها سازگار هستند؟
\end{enumerate}

\mnote{برای اشکال‌زدایی، می‌توانید از \lr{\textbackslash listfiles} در ابتدای سند استفاده کنید تا فهرست همه‌ی فایل‌های بارگذاری‌شده را ببینید.}

\section{خطاهای رایج}
\label{sec:modules-mistakes}

\subsection{بارگذاری ماژول در فاز اشتباه}
\label{sec:modules-mistake-phase}

\textbf{مشکل:} اگر ماژولی که به
\lr{xepersian} نیاز دارد را در فاز ۱
بارگذاری کنید، خطا می‌دهد.

\textbf{راه‌حل:} ماژول‌های وابسته به
\lr{xepersian} را در فاز ۳ بارگذاری کنید.

\subsection{بارگذاری بسته بعد از \lr{xepersian}}
\label{sec:modules-mistake-after}

\textbf{مشکل:} اگر بسته‌ای را بعد از
\lr{xepersian} بارگذاری کنید، تعاریف
\lr{bidi} بازنویسی می‌شوند و خطا می‌دهید.

\textbf{راه‌حل:} \lr{xepersian} را آخرین
بسته‌ی فاز ۲ نگه دارید.

\subsection{فراموش کردن \cmd{RequirePackage}}
\label{sec:modules-mistake-require}

\textbf{مشکل:} اگر ماژول را در
\file{matinbook.cls} بارگذاری نکنید،
استفاده نمی‌شود.

\textbf{راه‌حل:} همیشه ماژول‌های جدید را
در \file{matinbook.cls} اضافه کنید.

\section{تمرین‌ها}
\label{sec:modules-exercises}

\begin{exercise}
    فهرست ۱۹ ماژول \lr{MatinBook} را بنویسید
    و نقش هر یک را توضیح دهید.
    \label{exc:modules-list}
\end{exercise}

\begin{exercise}
    یک ماژول ساده بسازید که یک دستور جدید
    تعریف کند و آن را در فاز ۱ بارگذاری
    کنید.
    \label{exc:modules-add}
\end{exercise}

\begin{exercise}
    یک ماژول بسازید که به \lr{xepersian}
    نیاز دارد و آن را در فاز ۳ بارگذاری
    کنید.
    \label{exc:modules-phase3}
\end{exercise}

\begin{exercise}
    ترتیب بارگذاری ماژول‌ها را عوض کنید و
    خطای آن را ببینید. سپس آن را اصلاح
    کنید.
    \label{exc:modules-order}
\end{exercise}

\section{خلاصه}
\label{sec:modules-summary}

در این فصل، با ساختار ماژولار \lr{MatinBook}
آشنا شدیم:

\begin{itemize}
    \item \lr{MatinBook} از ۱۹ ماژول مستقل
    تشکیل شده که در سه دسته تقسیم می‌شوند:
    هسته، کاربردی، تم و لوکال.

    \item ماژول‌ها در سه فاز بارگذاری
    می‌شوند: قبل از \lr{xepersian}،
    \lr{xepersian}، بعد از \lr{xepersian}.

    \item \lr{xepersian} باید آخرین بسته‌ی
    فاز ۲ باشد.

    \item رنگ‌ها باید قبل از مصرف‌کنندگان
    بارگذاری شوند.

    \item فونت‌ها باید بعد از \lr{xepersian}
    تنظیم شوند.

    \item ماژول جدید با ساختار استاندارد
    (\cmd{NeedsTeXFormat}،
    \cmd{ProvidesPackage}، \cmd{PackageInfo}،
    \cmd{endinput}) ساخته می‌شود.

    \item خطاهای رایج شامل بارگذاری ماژول
    در فاز اشتباه، بارگذاری بسته بعد از
    \lr{xepersian}، و فراموش کردن
    \cmd{RequirePackage} است.
\end{itemize}

در فصل بعدی (\cref{chap:custom-theme})،
با ساخت تم سفارشی آشنا می‌شوید و یاد
می‌گیرید که چگونه یک تم کاملاً جدید طراحی
کنید.